% Artículo 1 — vigilancia poscomercialización nacional de VA-MENGOC-BC, 2017-2025 (español).
% Revista destino: Drug Safety (versión en español para difusión institucional y revistas hispanohablantes).
\documentclass[11pt]{article}
\newcommand{\DocLang}{es}
\input{papers/shared/preamble.tex}
\input{papers/shared/authors.tex}
\input{papers/shared/pending.tex}
\input{papers/p1-pharmacovigilance/shared/flow.tex}
\externaldocument{supplement}
\usepackage{lineno}
\onehalfspacing

\hypersetup{pdftitle={Seguridad poscomercializacion de la vacuna VA-MENGOC-BC en Cuba, 2017-2025},
  pdfauthor={Richard Alejandro Matos Arderi; Abel Ponce Gonzalez}}

\begin{document}

\begin{titlepage}
\SyntheticNotice
\begin{center}
{\LARGE\bfseries Seguridad poscomercialización de la vacuna antimeningocócica B y~C VA-MENGOC-BC
en Cuba, \VStudyFirstYear--\VStudyLastYear: análisis nacional de notificaciones espontáneas con
métodos de desproporcionalidad, clases latentes y aprendizaje automático\par}
\bigskip
\AuthorBlock
\end{center}
\bigskip
\noindent\textbf{Título abreviado:} Seguridad poscomercialización de VA-MENGOC-BC en Cuba\par
\noindent\textbf{Palabras clave:} eventos supuestamente atribuibles a la vacunación o
inmunización; farmacovigilancia; análisis de desproporcionalidad; vacunas antimeningocócicas;
vesículas de membrana externa; análisis de clases latentes; aprendizaje automático; Cuba\par
\bigskip
\noindent\fbox{\parbox{\dimexpr\linewidth-2\fboxsep-2\fboxrule}{%
\textbf{Puntos clave}
\begin{itemize}[leftmargin=*,itemsep=2pt,topsep=2pt]
\item De \PONReports{} notificaciones sin duplicados de eventos supuestamente atribuibles a la
  vacunación o inmunización recibidas por el sistema nacional de vigilancia de Cuba en
  \VStudyFirstYear--\VStudyLastYear, \PONTarget{} (\POPctTarget) incluyeron VA-MENGOC-BC, con una
  tasa global de notificación de \POPooledRate{} por \num{100000} dosis, \PORateVsExpected{} la
  tasa de referencia del programa para todas las vacunas infantiles.
\item Con un criterio bayesiano preespecificado, \POSignalsPrimaryAllOtherVaccines{} de
  \PONEventsScreened{} tipos de evento evaluados se notificaron de forma desproporcionada con
  VA-MENGOC-BC (\POSignalList); \PONSignalsRobustPhrase{} se mantuvo en los \PONDesigns{} diseños
  de comparación.
\item Los eventos notificados conjuntamente formaron \POLCABestK{} fenotipos latentes; la
  hospitalización se predijo con una discriminación \PODiscriminationWord{} en los años
  posteriores, por lo que los modelos basados en notificaciones deben describir y no clasificar
  casos.
\end{itemize}}}
\end{titlepage}

\linenumbers

\section*{Resumen}
\noindent\textbf{Introducción:} VA-MENGOC-BC, vacuna de vesículas de membrana externa contra la
enfermedad meningocócica por los serogrupos~B y~C, forma parte del esquema infantil cubano desde
1991, pero su seguridad poscomercialización no se ha analizado a escala nacional con métodos
actuales de farmacovigilancia.\par
\noindent\textbf{Objetivo:} Describir las tasas de notificación, la desproporcionalidad y la
estructura de los eventos supuestamente atribuibles a la vacunación o inmunización (ESAVI)
notificados con VA-MENGOC-BC en \VStudyFirstYear--\VStudyLastYear, así como los factores asociados
a la hospitalización.\par
\noindent\textbf{Métodos:} Análisis de desproporcionalidad de notificaciones individuales de casos
en la base de datos nacional de vigilancia de ESAVI de Cuba (acceso institucional), que abarca
\VStudyFirstYear--\VStudyLastYear{} y todas las provincias. Las notificaciones de \VNFiles{}
ficheros anuales se armonizaron, seudonimizaron y depuraron de duplicados. Las tasas de
notificación usaron como denominador las dosis administradas. VA-MENGOC-BC se comparó con el resto
de las vacunas mediante la razón de odds de notificación, la razón de notificación proporcional, el
componente de información bayesiano (IC) y la media geométrica bayesiana empírica, con un criterio
de señal preespecificado (límite inferior de credibilidad del 95\,\% IC\textsubscript{025}${}>0$
con al menos \PODispMinReports{} notificaciones) y tres comparadores de sensibilidad; no se realizó
un análisis caso a caso. El análisis de clases latentes resumió los eventos notificados
conjuntamente; la regresión logística de Firth y la potenciación del gradiente con validación
temporal describieron la hospitalización. El informe sigue READUS-PV.\par
\noindent\textbf{Resultados:} De \PONReports{} notificaciones, \PONTarget{} (\POPctTarget)
incluyeron VA-MENGOC-BC. La tasa global de notificación fue de \POPooledRate{} por \num{100000}
dosis (\CI{} \POPooledRateCI) y \POTrendWord{} en el tiempo (razón de tasas anual \POTrendRR;
\CI{} \POTrendRRCI). \POSignalsPrimaryAllOtherVaccines{} de \PONEventsScreened{} tipos de evento
cumplieron el criterio de señal (\POSignalList), \PONSignalsRobust{} de ellos con todos los
comparadores. Un modelo de \POLCABestK{} clases describió la notificación conjunta (índice de Rand
ajustado por remuestreo \POLCAAri). El mejor modelo de hospitalización (\POBestModel) alcanzó un
área bajo la curva ROC de \POBestAuroc{} en los años de prueba.\par
\noindent\textbf{Conclusiones:} Este análisis nacional y reproducible describe el perfil de
notificación poscomercialización de VA-MENGOC-BC. La desproporcionalidad genera hipótesis y no mide
riesgos: las señales que persisten con todos los comparadores requieren revisión clínica y
evaluación de causalidad, y el marco con control de divulgación permite una vigilancia anual
continua.

\clearpage
\section{Introducción}\label{sec:intro}

La enfermedad meningocócica invasiva sigue causando muertes y secuelas permanentes en el mundo, con
más de medio millón de casos y unas \num{85000} defunciones estimadas cada año
\cite{SierraGonzalez2019}. En Cuba, la incidencia alcanzó 5,9 casos por \num{100000} habitantes en
1980, cerca del 80\,\% causados por el serogrupo~B, y la enfermedad meningocócica fue declarada el
principal problema de salud del país \cite{SierraGonzalez2019}. VA-MENGOC-BC, desarrollada en el
Instituto Finlay, combina vesículas de membrana externa (VME, proteoliposomas) de una cepa del
serogrupo~B con el polisacárido capsular del serogrupo~C, adsorbidos en hidróxido de aluminio
\cite{Sotolongo2007,Uli2006,Masforrol2017}. Su eficacia del 83\,\% se demostró en un ensayo
aleatorizado, a doble ciegas y controlado con placebo \cite{Sierra1991}; tras la vacunación masiva
la vacuna se incorporó al programa nacional de inmunización, y su impacto acumulado sobre la
enfermedad por serogrupo~B superó el 95\,\% \cite{Rodriguez1999,SierraGonzalez2019}. Junto con las
experiencias de Noruega y Nueva Zelanda, el programa cubano es uno de los pocos en que vacunas de VME
de tipo silvestre controlaron epidemias por serogrupo~B \cite{Holst2009,Holst2013}.

Las vacunas de VME tienen un perfil de reactogenicidad característico. Las reacciones locales y la
fiebre son frecuentes en niños pequeños, mientras que los eventos graves y neurológicos son muy
raros \cite{Nokleby2007,Holst2009}. La pirogenicidad de las VME se ha relacionado con su endotoxina
unida a la membrana, tanto en modelos experimentales de la vacuna de cuatro componentes contra el
serogrupo~B \cite{Sheerin2019} como en estudios de formulación, en los que la adsorción en hidróxido
de aluminio redujo la pirogenicidad \cite{Rosenqvist1998}. Como VA-MENGOC-BC se administra a
lactantes junto con la vacuna pentavalente de células enteras de \emph{Bordetella pertussis}, que es
el producto más reactogénico del esquema cubano \cite{Galindo2012,CruzRodriguez2021}, atribuir la
fiebre y las reacciones locales a un solo producto no es sencillo.

Cuba estableció un sistema nacional de vigilancia de ESAVI en 1999 \cite{Galindo2012,Reed2007}. Su
primera década de datos (1999--2008) mostró una tasa global de 57,8 ESAVI por \num{100000} dosis
\cite{Galindo2012}, y el programa utiliza 50 por \num{100000} dosis como tasa nacional esperada
\cite{CruzRodriguez2021}. Desde entonces, los análisis publicados han sido descriptivos y
provinciales \cite{CruzRodriguez2021}, y ningún estudio ha examinado VA-MENGOC-BC con los métodos
hoy habituales en farmacovigilancia: análisis de desproporcionalidad con diseños de sensibilidad al
comparador \cite{Evans2001,vanPuijenbroek2002,Bate1998,Noren2013,DuMouchel1999}, tratamiento
explícito del enmascaramiento y de la coadministración \cite{Maignen2014,Salvo2013} e informe
transparente \cite{Fusaroli2024Statement,Fusaroli2024Elaboration}. El aprendizaje automático también
ha llegado a la farmacovigilancia, pero su valor añadido depende de la validación temporal y de la
calibración, no de la exactitud aparente \cite{Salas2022,Collins2024,VanCalster2019}.

Por ello analizamos las notificaciones espontáneas nacionales de \VStudyFirstYear--\VStudyLastYear{}
con cuatro objetivos: (i) estimar las tasas de notificación de ESAVI con VA-MENGOC-BC y su tendencia;
(ii) buscar notificación desproporcionada de eventos concretos frente al resto de las vacunas,
comprobando la robustez de cada señal ante la elección del comparador; (iii) describir la estructura
latente de los eventos notificados conjuntamente, y (iv) describir los factores asociados a la
hospitalización entre las notificaciones.

\section{Métodos}\label{sec:methods}

\subsection{Diseño, ámbito e informe}\label{sec:design}
Se realizó un análisis de desproporcionalidad de notificaciones individuales de casos,
complementado con análisis de tasas de notificación, de clases latentes y de predicción, en la base
de datos nacional de vigilancia de ESAVI de Cuba. La base recibe los formularios de investigación de
casos de los puntos de vacunación de todas las provincias a través de las unidades municipales y
provinciales de epidemiología \cite{Galindo2012,WHO2016AEFI}; su custodio es \VCustodian{} y los
datos se extrajeron el \VExtractionDate. El informe sigue la guía READUS-PV para análisis de
desproporcionalidad \cite{Fusaroli2024Statement,Fusaroli2024Elaboration} (lista de verificación en
el material suplementario), la declaración RECORD-PE para datos recogidos de forma rutinaria
\cite{Langan2018} y, para los modelos de hospitalización, TRIPOD+AI \cite{Collins2024}.

\subsection{Fuentes de datos y procesamiento}\label{sec:data}
El estudio utilizó \VNFiles{} ficheros anuales de notificaciones de ESAVI
(\VStudyFirstYear--\VStudyLastYear) que recogen, para cada notificación, datos demográficos, vacunas,
dosis, lotes y fabricantes, fechas de vacunación, de inicio y de notificación, \PONEventsScreened{}
indicadores de eventos predefinidos, antecedentes familiares, hospitalización y evolución. Las dosis
anuales de VA-MENGOC-BC administradas procedieron de las series del programa nacional de
inmunización, y la incidencia anual de enfermedad meningocócica (\POITSFirstYear--\POITSLastYear),
de las estadísticas nacionales. La base registraba \PONVaccines{} vacunas diferentes; las vacunas se
armonizaron con un diccionario versionado de nombres nacionales de productos y los eventos se
codificaron con los indicadores predefinidos del formulario nacional (el sistema no utiliza la
codificación MedDRA). Todo el procesamiento se ejecutó en una estación de trabajo institucional con
un flujo de análisis de código abierto (Python 3.11
\cite{Harris2020,Virtanen2020,McKinney2010,Seabold2010,Pedregosa2011,Hunter2007}; versión
\VCodeVersion, semilla aleatoria \VSeed) que registra la procedencia de cada ejecución
\cite{Sandve2013,MHRA2018}.

Como la estructura de los ficheros cambió entre años, los nombres de las columnas se asignaron a un
esquema versionado mediante alias, con informes de deriva del esquema, y cada valor se interpretó con
un estado explícito (válido, recuperado, inválido, ausente) para evaluar la calidad de los datos
según el marco armonizado de Kahn et al.\ \cite{Kahn2016}. \VFormSentence. Los nombres y las direcciones se
sustituyeron en la ingesta por resúmenes con clave (HMAC-SHA-256) cuya clave secreta se guarda fuera
del repositorio \cite{ENISA2019}; las fechas de nacimiento se convirtieron en edad, y los campos de
texto libre se procesaron localmente y nunca se publicaron, conforme a la ley cubana de protección
de datos personales \cite{Ley149}. Se eliminaron los registros duplicados exactos
(\VExactDuplicates) y se fusionaron los registros con la misma persona seudonimizada, fecha de
vacunación, vacunas y lotes (\VKeyDuplicatesMerged{} filas), con lo que quedaron \VReportsDedup{}
notificaciones. El año de análisis fue el de vacunación; la ventana de estudio conservó
\VReportsInWindow{} notificaciones (se excluyeron \VReportsOutsideWindow{} fuera de la ventana;
\figref{fig:flow}).

\begin{figure}[tbp]
\centering
\ReportFlow
\caption{Flujo de notificaciones de ESAVI desde los ficheros anuales hasta la población de
análisis. Las cifras proceden de la publicación con control de divulgación.}
\label{fig:flow}
\end{figure}

\subsection{Definición de la exposición y de los eventos}\label{sec:events}
Una notificación se atribuyó a VA-MENGOC-BC cuando la vacuna figuraba entre las administradas,
incluida la coadministración con otras vacunas. Los eventos fueron los \PONEventsScreened{}
indicadores del formulario nacional, que siguen la clasificación de la OMS en eventos comunes, raros
y graves \cite{WHO2016AEFI,CIOMS2012}; cuando existen, se citan como referencia las definiciones
correspondientes de la Brighton Collaboration
\cite{Marcy2004,Bonhoeffer2004Crying,Bonhoeffer2004Seizure,Ruggeberg2007,Kohl2007,Buettcher2007,Wise2007,Gidudu2008,Sejvar2007,Tapiainen2007},
aunque las notificaciones no se reclasificaron según ellas \cite{Kohl2008}. La evaluación de
causalidad del comité nacional \cite{WHO2018Causality} no se usó como criterio de inclusión y en este
estudio no se revisaron las notificaciones caso a caso. Antes del análisis, la fiebre de al menos 39
y 40\,\textdegree{}C, las reacciones locales severas y el absceso estéril se clasificaron como
reacciones esperadas de las vacunas de VME adsorbidas en aluminio
\cite{Nokleby2007,Holst2009,Jefferson2004}; las señales de cualquier otro evento se clasificaron como
emergentes.

\subsection{Tasas de notificación}\label{sec:m-rates}
Las tasas anuales de notificación se calcularon por \num{100000} dosis de VA-MENGOC-BC con
intervalos de Poisson exactos (Garwood) \cite{Garwood1936}. La tendencia lineal se estimó con un
modelo de cuasi-Poisson con el logaritmo de las dosis como desplazamiento (razón de tasas anual,
dispersión estimada con el estadístico de Pearson) y se comprobó con un modelo binomial negativo.
Las tasas se compararon de forma descriptiva con la tasa esperada del programa, de
\POExpectedRate{} por \num{100000} dosis, que se refiere a todas las vacunas infantiles
\cite{CruzRodriguez2021,Mahaux2016}.

\subsection{Análisis de desproporcionalidad}\label{sec:m-disp}
Para cada evento, una tabla de dos por dos cruzó las notificaciones con y sin VA-MENGOC-BC con las
notificaciones con y sin el evento. Se calcularon la razón de odds de notificación (ROR) con
intervalos de Woolf y corrección de Haldane cuando fue necesaria
\cite{vanPuijenbroek2002,Woolf1955,Haldane1956}, la razón de notificación proporcional (PRR) con la
$\chi^2$ de Yates \cite{Evans2001}, el componente de información (IC) de la red neuronal bayesiana de
propagación de la confianza con su límite inferior de credibilidad del 95\,\%,
IC\textsubscript{025} \cite{Bate1998,Noren2013}, y la media geométrica bayesiana empírica (EBGM) del
contractor gamma-Poisson multi-ítem, cuya distribución previa de dos gammas se estimó por máxima
verosimilitud sobre todos los pares vacuna--evento de la base \cite{DuMouchel1999,Szarfman2002}. El
criterio principal de señal, preespecificado, fue IC\textsubscript{025}${}>0$ con al menos
\PODispMinReports{} notificaciones; el número de métodos que cumplieron sus umbrales convencionales
(límite inferior de la ROR ${}>1$; PRR ${}\geq2$ con $\chi^2\geq4$; límite inferior de la EBGM
${}\geq2$) se presenta como medida de concordancia.

Todos los comparadores fueron vacunas, por lo que las medidas expresan la notificación relativa a
otras vacunas del mismo sistema y no el riesgo frente a la no vacunación. Se preespecificaron cuatro
diseños de comparación: el diseño principal (notificaciones con VA-MENGOC-BC frente al resto); un
diseño de comparador activo en lactantes, restringido a notificaciones en menores de
\POInfantMonths{} meses; un diseño restringido a notificaciones de una sola vacuna (VA-MENGOC-BC
administrada sola frente a otras vacunas administradas solas), y un diseño que elimina del grupo de
comparación las notificaciones que mencionan la vacuna pentavalente, el producto reactogénico
dominante, para evaluar el enmascaramiento \cite{Maignen2014}. El IC acumulado se recalculó año a
año para mostrar la estabilidad temporal de las estimaciones. Dado el carácter exploratorio de la
detección de señales, no se ajustó por comparaciones múltiples \cite{Rothman1990}; las señales son
hipótesis para la revisión clínica y no riesgos confirmados
\cite{Fusaroli2026Pitfalls,Fusaroli2026Causal}.

\subsection{Análisis de clases latentes de los eventos notificados conjuntamente}\label{sec:m-lca}
Los indicadores de eventos con una prevalencia de al menos el 0,5\,\% se modelaron con un análisis
de clases latentes (mezcla de distribuciones de Bernoulli independientes) ajustado con el algoritmo
de esperanza--maximización \cite{Dempster1977}, con una distribución previa Beta(1,5;\,1,5) sobre las
probabilidades de los ítems y \POLCAStarts{} inicios aleatorios por modelo. Los modelos de
\POLCAKMin{} a \POLCAKMax{} clases se compararon con el criterio de información bayesiano
\cite{Schwarz1978}; se informaron además la verosimilitud completada integrada \cite{Biernacki2000}
y la entropía \cite{Nylund2007}. La estabilidad se midió con la mediana del índice de Rand ajustado
\cite{Hubert1985} entre la clasificación original y \POLCABootstrap{} reajustes por remuestreo, tras
el emparejamiento óptimo de etiquetas \cite{Kuhn1955}. La pertenencia a las clases se comparó entre
las notificaciones con VA-MENGOC-BC y el resto mediante una prueba de $\chi^2$.

\subsection{Factores asociados a la hospitalización}\label{sec:m-hosp}
La hospitalización se modeló a partir de la edad, el sexo, la coadministración, la región, las
vacunas y los indicadores de eventos. Como los eventos raros y graves separan el desenlace, el modelo
logístico multivariable se ajustó por la verosimilitud penalizada de Firth \cite{Firth1993,Heinze2002},
tras eliminar los términos exactamente colineales, y sus razones de odds ajustadas se presentan con
intervalos de Wald. Un modelo de árboles con potenciación del gradiente (LightGBM) \cite{Ke2017}
recogió no linealidades e interacciones. Ambos modelos se entrenaron con \POTrainYears{} y se
evaluaron con \POTestYears{} (validación temporal), con el área bajo las curvas ROC y de
precisión--exhaustividad, la puntuación de Brier y el intercepto y la pendiente de calibración
\cite{VanCalster2019,Collins2024}; las predicciones logísticas usaron la corrección del intercepto
FLIC \cite{Puhr2017}. Las contribuciones exactas TreeSHAP describieron el modelo de potenciación
\cite{Lundberg2020}. Los modelos son explicativos y no se diseñaron para el triaje clínico.

\subsection{Patrones temporales y contexto epidemiológico}\label{sec:m-time}
Los recuentos mensuales de notificaciones se estabilizaron en varianza con la transformación de
Anscombe \cite{Anscombe1948} y se segmentaron con un algoritmo PELT exacto
\cite{Killick2012,Truong2020} (segmento mínimo de seis meses, penalización de tipo BIC). Para situar
la vacuna en su contexto, los casos nacionales anuales de enfermedad meningocócica
(\POITSFirstYear--\POITSLastYear) se analizaron como serie temporal interrumpida con un modelo
binomial negativo segmentado, tratando \POITSTransition{} (vacunación masiva e incorporación al
programa) como período de transición \cite{Bernal2017}.

\subsection{Control de divulgación, reproducibilidad y ética}\label{sec:m-sdc}
Solo salieron del entorno controlado resultados agregados. Se suprimieron los recuentos no nulos
menores que \VMinCell, se eliminaron los estadísticos derivados de recuentos suprimidos y una
supresión complementaria protegió las tablas de contingencia \cite{Hundepool2012}; un verificador
automático revisó cada fichero publicado. Cada número de este artículo se genera a partir de las
tablas publicadas, y el flujo de análisis, los datos sintéticos de prueba y la publicación están
disponibles abiertamente \cite{Wilkinson2016}. El análisis utilizó registros de vigilancia rutinaria
seudonimizados, tratados según el procedimiento institucional de gobierno de datos del Instituto
Finlay de Vacunas y la Ley 149/2022 \cite{Ley149}; la aprobación institucional se detalla en las
Declaraciones.

\section{Resultados}\label{sec:results}

\subsection{Notificaciones y características}\label{sec:r-reports}
El análisis incluyó \PONReports{} notificaciones, de las cuales \PONTarget{} (\POPctTarget)
incluyeron VA-MENGOC-BC y \PONTargetOnly{} de ellas no tuvieron vacunas coadministradas
(\tabref{tab:characteristics}). La mediana del tiempo entre la vacunación y la notificación fue de
\PODelayMedian{} días.

\begin{table}[tbp]
\centering\small
\caption{Características de las notificaciones de ESAVI con VA-MENGOC-BC y con otras vacunas. Los
valores son $n$ (\%) salvo indicación en contrario; se suprimen los recuentos menores que
\VMinCell, y [c] marca la supresión complementaria que impide recuperarlos por diferencia.}
\label{tab:characteristics}
\POTableCharacteristics
\end{table}

\subsection{Tasas de notificación}\label{sec:r-rates}
La tasa global de notificación de VA-MENGOC-BC entre \PORateFirstYear{} y \PORateLastYear{} fue de
\POPooledRate{} por \num{100000} dosis (\CI{} \POPooledRateCI), \PORateVsExpected{} la tasa de
referencia del programa, de \POExpectedRate{} por \num{100000} dosis para todas las vacunas
infantiles (\figref{fig:rates}; \tabref{tab:S-rates} del material suplementario). La tasa
\POTrendWord{} durante el período (razón de tasas anual \POTrendRR, \CI{} \POTrendRRCI;
$p$\POTrendP; dispersión de cuasi-Poisson \PODispersion).

\begin{figure}[tbp]
\centering
\includegraphics[width=\textwidth]{p1_rates}
\caption{Notificación en el tiempo. (a) Número mensual de notificaciones de ESAVI para todas las
vacunas y para VA-MENGOC-BC; las líneas verticales marcan los puntos de cambio detectados por PELT en
la serie de todas las vacunas. (b) Tasa anual de notificación de VA-MENGOC-BC por \num{100000} dosis
con intervalo de confianza exacto del 95\,\% (banda); la línea horizontal es la tasa esperada del
programa para todas las vacunas infantiles.}
\label{fig:rates}
\end{figure}

\subsection{Desproporcionalidad}\label{sec:r-disp}
Con el comparador principal, \POSignalsPrimaryAllOtherVaccines{} de los \PONEventsScreened{} tipos
de evento evaluados cumplieron el criterio de señal preespecificado (\POSignalList;
\tabref{tab:disproportionality}, \figref{fig:forest}): \PONSignalsExpectedPhrase{} de reacciones
esperadas (\POSignalExpectedList) y \PONSignalsEmergingPhrase{} de eventos no clasificados como
esperados (\POSignalEmergingList). El número de señales fue \POSignalsInfantActiveComparator{} en el
diseño de comparador activo en lactantes, \POSignalsExcludingCoadministration{} en las
notificaciones de una sola vacuna y \POSignalsExcludingPentavalentMasking{} al eliminar del
comparador las notificaciones con vacuna pentavalente (\tabref{tab:designs};
Tablas~\ref{tab:S-disp-primary}--\ref{tab:S-disp-masking} del material suplementario);
\PONSignalsRobustPhrase{} se mantuvo en los \PONDesigns{} diseños (\POSignalRobustList). El IC
acumulado de cada evento se muestra año a año en la \tabref{tab:S-cumic} del material
suplementario. Para la fiebre de al menos 39\,\textdegree{}C, la ROR fue \POEvFeverThreeNineROR{}
(\CI{} \POEvFeverThreeNineRORCI) y el IC\textsubscript{025}, \POEvFeverThreeNineICLow; para las
reacciones locales severas fueron \POEvSevereLocalReactionROR{} (\CI{}
\POEvSevereLocalReactionRORCI) y \POEvSevereLocalReactionICLow.

\begin{table}[tbp]
\centering\small
\caption{Desproporcionalidad de los eventos notificados con VA-MENGOC-BC frente al resto de las
vacunas (diseño principal), para los eventos con al menos \VMinCell{} notificaciones. $a$,
notificaciones con VA-MENGOC-BC y el evento; $E$, valor esperado de $a$; IC\textsubscript{025},
límite inferior de credibilidad del 95\,\% del componente de información; EB\textsubscript{05},
percentil 5 inferior de la media geométrica bayesiana empírica; Métodos, número de los cuatro métodos
que cumplen sus umbrales; Señal, criterio preespecificado (IC\textsubscript{025}${}>0$,
$a\geq\PODispMinReports$).}
\label{tab:disproportionality}
\POTableDispMain
\end{table}

\begin{table}[tbp]
\centering\small
\caption{Número de señales de desproporcionalidad en los cuatro diseños de comparación
preespecificados.}
\label{tab:designs}
\POTableDesigns
\end{table}

\begin{figure}[tbp]
\centering
\includegraphics{p1_forest}
\caption{Razones de odds de notificación (IC 95\,\%, escala logarítmica) de los eventos notificados
con VA-MENGOC-BC frente al resto de las vacunas (diseño principal) y en el diseño de comparador
activo en lactantes. No se muestran los eventos con menos de \VMinCell{} notificaciones en el diseño
principal; los cuatro diseños se presentan en las
Tablas~\ref{tab:S-disp-primary}--\ref{tab:S-disp-masking} del material suplementario.}
\label{fig:forest}
\end{figure}

\subsection{Fenotipos latentes de los eventos notificados conjuntamente}\label{sec:r-lca}
El criterio de información bayesiano seleccionó un modelo de \POLCABestK{} clases
(\tabref{tab:S-lca} del material suplementario), cuya clasificación fue estable entre muestras de
remuestreo (mediana del índice de Rand ajustado \POLCAAri). La pertenencia a las clases
\POLCAGroupWord{} entre las notificaciones con VA-MENGOC-BC y el resto ($\chi^2=\POLCAChi$,
\POLCADf{} grados de libertad; $p$\POLCAP; \figref{fig:lca}).

\begin{figure}[tbp]
\centering
\includegraphics[width=\textwidth]{p1_lca}
\caption{Clases latentes de los eventos notificados conjuntamente. (a) Probabilidad de cada evento
dentro de cada clase latente (se imprimen los valores de al menos 0,2). (b) Proporción de
notificaciones asignadas a cada clase entre las notificaciones con VA-MENGOC-BC y con otras
vacunas.}
\label{fig:lca}
\end{figure}

\subsection{Factores asociados a la hospitalización}\label{sec:r-hosp}
El período de entrenamiento incluyó \PONTrain{} notificaciones (\POEventsTrain{} hospitalizaciones)
y el de prueba, \PONTest{} notificaciones (\POEventsTest{} hospitalizaciones). En los años de
prueba, el modelo de \POBestModel{} tuvo la mayor discriminación (área bajo la curva ROC
\POBestAuroc; \tabref{tab:models}), que calificamos como \PODiscriminationWord. Una pendiente de
calibración menor que uno indica predicciones demasiado extremas en los años de prueba, y una
pendiente mayor que uno, predicciones demasiado moderadas. Las razones de odds ajustadas del modelo
de Firth se presentan en la \tabref{tab:S-logistic} del material suplementario y la importancia
TreeSHAP, en la \figref{fig:shap}.

\begin{table}[tbp]
\centering\small
\caption{Validación temporal de los modelos de hospitalización (entrenados con \POTrainYears,
evaluados con \POTestYears). AUROC y AUPRC, áreas bajo las curvas característica operativa del
receptor y de precisión--exhaustividad.}
\label{tab:models}
\POTableModels
\end{table}

\begin{figure}[tbp]
\centering
\includegraphics{p1_shap}
\caption{Contribución TreeSHAP absoluta media de los doce predictores más influyentes al logaritmo
de odds predicho de hospitalización en los años de prueba (modelo de potenciación del gradiente).
Los contrastes por región se refieren a la región occidental.}
\label{fig:shap}
\end{figure}

\subsection{Patrones temporales y contexto}\label{sec:r-time}
La segmentación PELT identificó \PONBreaksAllPhrase{} en el número mensual de notificaciones
(\POBreaksAll). En la serie nacional de enfermedad meningocócica, el modelo segmentado mostró
\POITSWord{} tras la transición de \POITSTransition{} (razón de tasas del cambio de nivel
\POITSLevelRR, \CI{} \POITSLevelRRCI); la incidencia media fue de \POITSRatePre{} por \num{100000}
habitantes antes de la campaña y de \POITSRatePost{} después (\figref{fig:S-its} del material
suplementario).

\section{Discusión}\label{sec:discussion}

\subsection{Hallazgos principales}\label{sec:d-principal}
En \VStudyFirstYear--\VStudyLastYear, VA-MENGOC-BC representó el \POPctTarget{} de las
notificaciones de ESAVI, con una tasa global de notificación de \POPooledRate{} por \num{100000}
dosis que \POTrendWord{} durante el período. El análisis de desproporcionalidad señaló
\POSignalList{} con el comparador principal, y \POSignalRobustList{} con todos los comparadores. Los
eventos notificados conjuntamente se organizaron en \POLCABestK{} fenotipos, y la información de las
notificaciones predijo la hospitalización con una discriminación \PODiscriminationWord.

\subsection{Interpretación}\label{sec:d-interpretation}
Los eventos notificados con más frecuencia con VA-MENGOC-BC fueron \POTopEvents. La fiebre y las
reacciones locales son las reacciones esperadas de las vacunas de VME adsorbidas en aluminio
\cite{Nokleby2007,Holst2009}, cuya endotoxina determina la pirogenicidad
\cite{Sheerin2019,Rosenqvist1998}. En una base que solo contiene vacunas, la desproporcionalidad
mide cuánto difiere el perfil de las notificaciones de VA-MENGOC-BC del de otras vacunas. Como la
mayoría de los comparadores son también reactogénicos y la vacuna pentavalente de células enteras se
coadministra con VA-MENGOC-BC, una señal que desaparece al eliminar la coadministración o la vacuna
pentavalente se explica mejor por la atribución y el enmascaramiento que por el producto
\cite{Maignen2014,Salvo2013}; por el contrario, las señales que persisten con todos los comparadores
son las que más merecen revisión clínica. El diseño estructurado con varios comparadores es, por
tanto, la principal salvaguarda de este análisis, junto con el criterio preespecificado y el informe
transparente de todos los eventos evaluados \cite{Fusaroli2024Elaboration,Fusaroli2026Pitfalls}.
Las señales de reacciones esperadas (\POSignalExpectedList) cuantifican cuánto más aparece en las
notificaciones la reactogenicidad conocida de una vacuna de VME que en las de otras vacunas; las
señales de eventos emergentes (\POSignalEmergingList) son las que requieren revisión caso a caso
antes de cualquier interpretación.

La tasa de notificación debe leerse según las características del sistema cubano: un formulario
obligatorio cumplimentado en el punto de vacunación, el seguimiento activo de los servicios de
epidemiología y una sólida cultura de notificación \cite{Galindo2012,Reed2007}. En estos sistemas,
las tasas de notificación reflejan tanto la sensibilidad de la vigilancia como la frecuencia de las
reacciones, como mostró la experiencia de Pinar del Río \cite{CruzRodriguez2021}, y la OMS utiliza
las tasas de notificación como indicador del desempeño de la vigilancia \cite{Lei2018}. La tasa de
referencia del programa se refiere a todas las vacunas, por lo que la comparación con VA-MENGOC-BC es
solo orientativa \cite{Mahaux2016}.

Las clases latentes ofrecen una alternativa basada en los datos a la evaluación evento a evento:
describen los síndromes tal como se notifican, separan las presentaciones febriles de las locales o
alérgicas y permiten comparar la distribución de fenotipos entre vacunas. Los modelos de
hospitalización alcanzaron una discriminación \PODiscriminationWord{} en los años posteriores.
\POHospInterpretation, por lo que los modelos basados en notificaciones no deben usarse para el
triaje de casos \cite{VanCalster2019,Collins2024}.

Los hallazgos corresponden a un sistema de vigilancia maduro y obligatorio, con alta cobertura de
vacunación, y a una vacuna de VME de tipo silvestre adsorbida en aluminio que se administra con una
vacuna pentavalente de células enteras; su extrapolación a otros contextos o formulaciones de VME
requiere el mismo enfoque de sensibilidad al comparador. Dos diseños adicionales pondrían a prueba
las señales: una serie de casos autocontrolada con registros vinculados de inmunización y
hospitalización, que elimina la confusión invariante en el tiempo, y el análisis de los atributos de
calidad de los lotes administrados, que aborda la variabilidad de la fabricación.

\subsection{Fortalezas y limitaciones}\label{sec:d-limitations}
El estudio abarca todas las notificaciones recibidas a escala nacional en
\VStudyFirstYear--\VStudyLastYear, armoniza ficheros cuya estructura cambió con el tiempo,
preespecifica su criterio de señal y sus diseños de comparación, y genera cada número a partir de
una publicación verificada y con control de divulgación, de modo que el análisis puede reproducirse
y ampliarse cada año.

Sus limitaciones son las de la notificación espontánea, reforzadas por rasgos propios de los datos.
En primer lugar, la infranotificación y la notificación estimulada afectan a los numeradores
\cite{Hazell2006,Pariente2007}; las dosis administradas, único denominador disponible, no estaban
desagregadas por edad, dosis ni provincia. En segundo lugar, los indicadores de eventos no se
validaron frente a las definiciones de Brighton \cite{Kohl2008}, y las descripciones en texto libre
se excluyeron del análisis para proteger la confidencialidad. En tercer lugar, las notificaciones
atribuidas a VA-MENGOC-BC incluyen vacunas coadministradas, y la desproporcionalidad no puede separar
sus contribuciones; los diseños de sensibilidad abordan esta limitación pero no la eliminan. En
cuarto lugar, todos los comparadores son vacunas, por lo que la desproporcionalidad no mide el riesgo
frente a la no vacunación, y una señal es una hipótesis para la evaluación de causalidad
\cite{WHO2018Causality,Fusaroli2026Causal}. En quinto lugar, la serie temporal interrumpida de la
incidencia nacional es ecológica y sirve solo de contexto. Por último, los recuentos pequeños se
suprimieron en la publicación, por lo que los eventos graves raros solo pueden describirse de forma
cualitativa.

\subsection{Implicaciones}\label{sec:d-implications}
El flujo de análisis desarrollado para este estudio puede ejecutarse cada año con los nuevos
ficheros y convertir un análisis retrospectivo en una vigilancia continua y auditable. Las señales
que persisten con todos los comparadores deben ser revisadas con las historias clínicas por el
comité nacional de causalidad, y la vinculación de las notificaciones con los atributos de calidad
de los lotes administrados permite estudiar si la variabilidad de la fabricación contribuye a la
reactogenicidad.

\section{Conclusiones}\label{sec:conclusions}
Este análisis nacional describe el perfil de notificación poscomercialización de VA-MENGOC-BC con
métodos transparentes en sus supuestos y reproducibles a partir de una publicación pública y con
control de divulgación. Las señales que persisten en todos los diseños de comparación definen las
prioridades de la revisión clínica, y el marco ofrece una base para la vigilancia continua de la
seguridad del programa cubano de inmunización.

\section*{Declaraciones}
\input{papers/p1-pharmacovigilance/shared/declarations_es.tex}

\bibliographystyle{vancouver}
\bibliography{papers/shared/bib/pubmed,papers/shared/bib/manual}

\end{document}
